(1) 行进光栅的两傅里叶谐波为 \(e^{\pm i(Kz-\Omega_at)}\)，故一阶衍射的选择规则是
\[
(\omega_1,k_{1z})=(\omega_0\pm\Omega_a,k_{0z}\pm K).
\]
选上转换 \(\omega_1=\omega_0+\Omega_a\)，准确Bragg匹配对应 \(\Delta=k_{1z}-k_{0z}-K=0\)；若光在晶体中的色散关系使 \(\Delta\ne0\)，沿传播长度的相对相位会失配。静态相位光栅不具有这样的必然频率偏移。

(2) 引入 \(a=Ae^{-i\Delta z/2},\ b=Be^{i\Delta z/2}\)，得
\[
\frac d{dz}\binom A B
=i\begin{pmatrix}\Delta/2&\kappa\\\kappa&-\Delta/2\end{pmatrix}\binom A B.
\]
矩阵平方是 \(\nu^2 I\)，其中 \(\nu=\sqrt{\kappa^2+\Delta^2/4}\)，所以
\(B(L)=i(\kappa/\nu)\sin(\nu L)\)。复相因子不改强度，因而
\[
\boxed{P(L,\Delta)=\frac{\kappa^2}{\kappa^2+\Delta^2/4}
\sin^2[L\sqrt{\kappa^2+\Delta^2/4}]}.
\]
此解是无衰减双模耦合方程的精确矩阵指数结果，不是弱衍射的Born近似。

(3) \(\Delta=0\) 时 \(P=\sin^2(\kappa L)\)，首次全转换长度
\(L_c=\pi/(2\kappa)\)。若 \(\Delta\ne0\)，
\(P\le\kappa^2/(\kappa^2+\Delta^2/4)<1\)。
固定长度且 \(\kappa L\ll1\)，得
\[
P=(\kappa L)^2\operatorname{sinc}^2(\Delta L/2)+O((\kappa L)^4),
\]
主阶首零点 \(|\Delta|=2\pi/L\)；精确有限耦合的零点可以略有偏移。
若 \(L=L_c\)，小失配展开为
\(P=1-\Delta^2/(4\kappa^2)+O(\Delta^4/\kappa^4)\)，即首个效率损失为二次阶。

(4) 将耦合方程与共轭方程相乘相加，可得
\(\frac d{dz}(|a|^2+|b|^2)=0\)，故其值始终为1。
若每秒入射光子数 \(\dot N\)，上转换级和未衍射级光功率分别是
\(\hbar(\omega_0+\Omega_a)\dot N P\) 与
\(\hbar\omega_0\dot N(1-P)\)；其和比入射光功率多
\(\hbar\Omega_a\dot N P\)，恰从声波泵浦提取。
光子通量守恒不等于光功率守恒；此双模模型要求声泵足够强、其他衍射级失谐以及光学损耗可忽略。
